\documentclass[11pt]{article}
\usepackage{series}
\usepackage{mathrsfs}

\title{World-in-World Genesis:\\
Local Comparison Geometry and a Globalization Program\\
\large Corrected fifth edition}
\author{Peter Nero}
\date{July 2026}

\begin{document}
\maketitle

\begin{abstract}
We give the precise mathematical content of the world-in-world proposal and
separate it from its physical interpretation.  The central object is a
comparison field $Q_{\rm WW}\in\Gamma(\operatorname{Hom}(TP,TI))$ between two
oriented rank-three bundles.  Its nine local components split, at a nonsingular
background, into three orientation and six strain directions.  A selected
flag resolves the strain sector into scalar, diagonal-shape, and shear sectors
of dimensions $1+2+3$.  Therefore
$1+3\times3=(1+3)+(1+2+3)=4+6$ is an exact component identity.  It is not
ordinary manifold-dimension multiplication and does not derive a
ten-dimensional spacetime, Lorentzian signature, gravity, matter, or
quantization.  We state the missing globalization theorem to the selected q79
Fu--Yau carrier, distinguish the shared phase circle from physical time, and
formulate the additional gates needed for each proposed physical emergence.
The resulting paper is a rigorous local geometry theorem plus a falsifiable
research program.
\end{abstract}

\section*{Revision note for this edition}
\addcontentsline{toc}{section}{Revision note for this edition}
\begin{description}
\item[Supersedes.] \emph{World-in-World Genesis: A Proto-Geometric Origin of
Time, Gravity, Matter, and Quantization}, version~4.
\item[Reason.] A three-dimensional base and three-dimensional fiber were
multiplied into nine manifold dimensions, and the resulting count was used to
claim automatic emergence of spacetime and several physical sectors.
\item[Resolution.] Version~5 replaces that construction by
$Q_{\rm WW}\in\Gamma(\operatorname{Hom}(TP,TI))$, proves the local polar and
$1+2+3$ strain decompositions, and states separate globalization/physical
gates.
\item[Retained result.] The exact component identity
$1+3\times3=(1+3)+(1+2+3)=4+6$ survives as local comparison geometry.
\item[Remaining boundary.] The q79 bundle-and-connection intertwiner,
Lorentzian completion, gravity, matter, and quantization are not yet generated
by this identity.
\end{description}

\tableofcontents

\section{Claim discipline}

This paper has three layers:
\begin{description}
\item[Proved local geometry.] Bundle dimensions, polar decomposition, and the
orthogonal component splittings stated below.
\item[Selected downstream data.] The q79 trace-split carrier and Fu--Yau branch
as recorded by their independent theorem packets.
\item[Research program.] The proposed identification of local strain sectors
with q79 vertical operators and, after that, with physical spacetime fields.
\end{description}

No statement in the third layer is promoted by suggestive equality of
dimensions.  A bridge must identify maps, transition functions, connections,
actions, and observables.

\section{The local comparison field}

Let $B$ be a smooth manifold and let $TP,TI\to B$ be oriented Euclidean
rank-three bundles.  The labels ``outer'' and ``inner'' describe the roles of
the two fibers; they do not create separate universes or coordinate factors.

\begin{definition}[World-in-world field]
A world-in-world field is
\[
 Q_{\rm WW}\in\Gamma\!\left(\operatorname{Hom}(TP,TI)\right).
\]
Under changes of oriented orthonormal frames
$g_P,g_I:B\to SO(3)$, its local matrix transforms as
\[
 Q_{\rm WW}\longmapsto g_I Q_{\rm WW}g_P^{-1}.
\]
\end{definition}

This transformation law is essential: the nine entries are not nine scalar
fields independent of frame choice.

\begin{theorem}[Bundle dimension versus comparison components]
If $\dim B=3$, the total space of either rank-three vector bundle has dimension
six.  The bundle $\operatorname{Hom}(TP,TI)$ has rank nine, so a section has
nine local fiber components.
\end{theorem}

\begin{proof}
Local bundle charts have the forms $U\times\mathbb R^3$ and
$U\times\operatorname{Mat}(3,\mathbb R)$, respectively.  Dimensions are
$3+3=6$ and ranks are $3\cdot3=9$.
\end{proof}

Thus the original intuition survives only as ``three-dimensional outer space
plus a $3\times3$ comparison field.''  It does not survive as three independent
three-dimensional manifolds.

\section{Polar and strain decomposition}

Assume $Q_{\rm WW}(b)\in GL^+(3,\mathbb R)$ on an admissible domain.  Polar
decomposition is unique:
\[
 Q_{\rm WW}=RU,
 \qquad R\in SO(3),
 \qquad U=(Q_{\rm WW}^{T}Q_{\rm WW})^{1/2}>0.
\]
The logarithmic strain is
\[
 S=\log U\in\operatorname{Sym}(3,\mathbb R).
\]

\begin{theorem}[Orientation--strain normal form]
At every nonsingular background, the tangent space of the comparison carrier
has the orthogonal Frobenius decomposition
\[
 \operatorname{Mat}(3,\mathbb R)
 =\mathfrak{so}(3)\oplus\operatorname{Sym}(3,\mathbb R),
 \qquad 9=3+6.
\]
\end{theorem}

\begin{proof}
Every matrix $X$ decomposes uniquely as
$X=(X-X^T)/2+(X+X^T)/2$.  Skew and symmetric matrices are orthogonal for
$\langle X,Y\rangle_F=\operatorname{tr}(X^TY)$ and have dimensions three and
six.
\end{proof}

Choose an orthonormal flag, equivalently a preferred ordered frame modulo the
declared residual symmetry.  For $S\in\operatorname{Sym}(3)$ define
\[
 P_1S=\frac{\operatorname{tr}S}{3}I_3,
 \qquad
 P_2S=\operatorname{diag}(S)-P_1S,
 \qquad
 P_3S=S-\operatorname{diag}(S).
\]

\begin{theorem}[Selected $1+2+3$ strain split]
The maps $P_1,P_2,P_3$ are mutually orthogonal projectors with ranks
$1,2,3$, and
\[
 S=P_1S+P_2S+P_3S,
 \qquad
 \|S\|_F^2=\sum_{j=1}^3\|P_jS\|_F^2.
\]
\end{theorem}

\begin{proof}
$P_1S$ is scalar diagonal, $P_2S$ is traceless diagonal, and $P_3S$ has zero
diagonal.  These subspaces are pairwise Frobenius-orthogonal, their dimensions
sum to six, and the displayed formulas follow directly.
\end{proof}

The $1+2+3$ split is flag-dependent.  Without a selected flag, the canonical
$SO(3)$-equivariant split is $1+5$ into trace and traceless symmetric tensors.
The flag is therefore source data that must either be selected by q79 geometry
or declared as an assumption.

\section{What the ten-component identity says}

If one additional scalar $u$ records an ordering coordinate, the local data
have the count
\[
 \underbrace{1}_{u}
 +\underbrace{9}_{Q_{\rm WW}}
 =\underbrace{(1+3)}_{u\;\text{and orientation}}
 +\underbrace{(1+2+3)}_{\text{strain}}
 =4+6.
\]

\begin{proposition}[No automatic spacetime theorem]
The preceding equality does not imply a splitting
$TM_{10}\cong TY_4\oplus TX_6$, does not select a Lorentzian metric on a
four-dimensional factor, and does not identify the strain bundle with a
compact internal tangent bundle.
\end{proposition}

\begin{proof}
A component count contains no transition functions for $TM_{10}$, no
integrable distributions, no metric of Lorentzian inertia, and no compactness
or topology data.  Each is independent structure.
\end{proof}

The identity is nevertheless useful: it gives a precise local carrier whose
four orientation/order and six strain components can be compared with a
separately supplied $4+6$ physical realization.

\section{Shared circle and physical time}

Let $L_{\rm shared}\to X_6$ be a Hermitian line bundle with unitary connection.
Its circle fibers encode common phase or holonomy.  Reusing this line in all
three rank lanes does not add three circle coordinates and does not create a
seventh internal dimension.

Physical time in the canonical completion is a noncompact ordering variable
on a Lorentzian base, represented by a propagator $U(t_2,t_1)$.  A phase map
\[
 \mathbb R\longrightarrow S^1,
 \qquad t\longmapsto e^{i\omega t},
\]
can synchronize phase with time after $t$ and $\omega$ are supplied, but it is
many-to-one and cannot make compact phase identical to physical time.  An
arrow of time further requires an asymmetric dynamical law or boundary
condition.

\section{Selected q79 global carrier}

For the selected degree-three cover $\pi:C\to B$, define
\[
 \mathcal A=\pi_*\mathcal O_C,
 \qquad
 \mathcal A_0=\ker\operatorname{Tr},
 \qquad
 \mathcal A=\mathcal O\oplus\mathcal A_0.
\]
The selected common carrier is
\[
 \mathcal H_{\rm q79}
 =L_{\rm shared}\otimes
 (\mathcal O\oplus\mathcal A_0\oplus\mathcal A),
 \qquad \operatorname{rank}=1+2+3=6.
\]

The trace split is global and intrinsic to the cover.  It is not a global
ordered-sheet flag: connected transitive monodromy prevents such an ordering.
The signed sheet representation has a local binary-dihedral spin lift, but
strict global Spin closure still requires all branch-complement relator signs
or the equivalent obstruction class.

The q79 Fu--Yau branch is the current selected compactification candidate.
The auxiliary space $L(3,1)\times\mathrm{Nil}_3$ is not homeomorphic to it and
is used only as an effective model or operator mnemonic.  Literal
$S^1\times\mathrm{Lens}\times\mathrm{Nil}$ and literal manifold nesting are
not premises of this paper.

\section{The globalization theorem}

The central open map is a bundle morphism
\[
 \mathfrak I:
 \mathbb RI_3\oplus\mathcal D_0\oplus\mathcal O
 \longrightarrow
 L_{\rm shared}\otimes
 (\mathcal O\oplus\mathcal A_0\oplus\mathcal A).
\]

\begin{conjecture}[Selected world-in-world/q79 intertwiner]
On the selected q79 Fu--Yau branch there is a source-selected flag and an
isometric bundle isomorphism $\mathfrak I$ preserving the $1+2+3$ filtration
such that
\[
 \mathfrak I\nabla^{\rm WW}=\nabla^{\rm HYM}\mathfrak I,
 \qquad
 \mathfrak I H_{\rm WW}=H_{\rm q79}\mathfrak I,
\]
and likewise for the retarded and overlap kernels used in the numerical
source packets.
\end{conjecture}

The conjecture has four independent tests:
\begin{enumerate}
\item transition functions agree on all overlaps;
\item the selected flag is monodromy-compatible rather than an ordered-sheet
choice;
\item the two inner products and connections are intertwined; and
\item the resulting operator coefficients equal the source coefficients used
downstream, with a declared normalization and error certificate.
\end{enumerate}

Rank matching proves none of these clauses, but it makes the conjecture
well-typed and finite.

\section{Physical emergence gates}

The original narrative can now be expressed as separate targets.

\subsection{Time}
A noncompact parameter, physical propagator, and arrow-producing condition are
required.  Bookkeeping order or phase winding alone is insufficient.

\subsection{Gravity}
A Lorentzian metric or tetrad, connection, action or equations, constraints,
and a stress tensor are required.  The local polar factor $R$ does not derive
the Einstein--Hilbert action.

\subsection{Matter}
A spin bundle, finite representation data, Dirac operator, chirality and real
structure, and a particle/observable interpretation are required.  The local
$1+2+3$ strain split alone does not select Standard Model representations.

\subsection{Quantization}
A complex state space, observable algebra, dynamics, probability functional,
and gauge/constraint treatment are required.  Discrete topology or a compact
resolvent can produce discrete spectra but not the Born rule by itself.

\subsection{Mass and scalar relaxation}
Mass requires a canonically normalized pole or dispersion relation.  A Higgs
mode requires a selected scalar representation and couplings.  Positivity of
a six-dimensional strain Hessian does not imply a unique radial Higgs.

\section{Relation to current calculations}

The finite $27\times27$ matrix, charged Yukawa profile, CKM, electroweak,
threshold, and precision packets close embedded renormalized-SM equivalence at
the declared one-shared-physical-primitive/profile standard.  They constitute
a nontrivial downstream realization.  Standard BRST/Faddeev--Popov
quantization and measured renormalized profile data are admitted at that
standard; strict no-knob selection remains an upgrade.  These computations do
not retroactively prove the world-in-world globalization conjecture.

\section{Conclusion}

World-in-world geometry has a clean surviving core.  It is a rank-three
comparison field with an exact orientation/strain decomposition and a selected
$1+2+3$ local normal form.  Its physical promise is concentrated into one
testable globalization problem rather than many automatic-emergence claims.
Failure to construct the q79 intertwiner would falsify the proposed unified
identification while leaving the local linear algebra intact.

\begin{thebibliography}{9}
\bibitem{Helgason}
S.~Helgason, \emph{Differential Geometry, Lie Groups, and Symmetric Spaces},
Academic Press, 1978.

\bibitem{FuYau}
J.-X.~Fu and S.-T.~Yau,
\emph{The theory of superstring with flux on non-Kahler manifolds and the
complex Monge--Ampere equation}, J. Differential Geom. 78 (2008).

\bibitem{MTTRecon}
P.~Nero, \emph{MTT Foundational Geometry Reconciliation}, internal theorem
and verification packet, 2026.
\end{thebibliography}

\end{document}
